This chapter describes the supervisor-facing interface and presents the
end-to-end validation results of the complete system.

% ═══════════════════════════════════════════════════════════
\section{Dashboard Architecture}
\label{sec:dashboard-arch}
% ═══════════════════════════════════════════════════════════

The Streamlit dashboard is the only component the branch supervisor
interacts with directly. Every other service --- RedPanda, Redis, the
scoring pipeline --- is invisible infrastructure. The dashboard must
therefore surface the right information clearly while remaining responsive
under continuous use.

\subsection{Why the Dashboard Cannot Consume Kafka Directly}

Streamlit's execution model re-runs the entire Python script from top to
bottom on every user interaction: clicking a button, changing a filter, or
navigating between pages triggers a full re-execution. A Kafka consumer
embedded in the script would be destroyed and recreated on each interaction,
losing its consumer group offset and replaying messages from the last
committed position. This makes direct Kafka consumption architecturally
incompatible with Streamlit.

The solution is the sink consumer pattern described in
Chapter~\ref{ch:streaming}: two independent consumers write streaming
results to PostgreSQL (the \texttt{alerts} table) and Redis (risk history)
continuously. The dashboard reads from PostgreSQL, which provides stable,
queryable, filterable access to all historical alerts without the
limitations of a streaming consumer.

\subsection{Data Access Pattern}

Reads come directly from PostgreSQL using \texttt{psycopg2} with
Streamlit's \texttt{@st.cache\_data} decorator for query-level caching.
Cache TTLs are tuned per page: 10 seconds for the alerts list (where
freshness matters) and 30 seconds for KPI aggregations (where stability
matters more than immediacy).

A \texttt{safe\_query()} wrapper in \texttt{src/dashboard/db.py} handles
connection resilience. If a query fails due to a stale connection --- common
after Docker restarts or idle timeouts --- the wrapper clears the cached
connection, establishes a new one, and retries the query transparently. This
single-retry pattern avoids the complexity of a connection pool while
handling the most common failure mode.

Writes go exclusively through FastAPI endpoints. The dashboard never
executes \texttt{INSERT} or \texttt{UPDATE} statements against the database.
This preserves the single-writer principle and ensures that all write-path
logic --- validation, audit logging, the feedback loop --- lives in one
place.

\subsection{Real-Time Refresh}

The \texttt{streamlit-autorefresh} component triggers a page re-execution
every 10 seconds. Streamlit does not support native WebSocket push from
server to client, so periodic polling is the pragmatic alternative. The
10-second interval balances responsiveness (a new BLOCK-tier alert appears
within 10 seconds of being written to PostgreSQL) against unnecessary
database load from overly aggressive polling.

\subsection{Navigation Structure}

The dashboard uses Streamlit's multi-page application structure, where each
page is an independent Python file in a \texttt{pages/} directory. Navigating
to a page executes only that page's script, avoiding the overhead of loading
data for all pages on every interaction. A shared \texttt{components.py}
module renders the sidebar on every page, displaying pipeline status
indicators and pending alert counts grouped by tier.


% ═══════════════════════════════════════════════════════════
\section{Alerts Page}
\label{sec:dashboard-alerts}
% ═══════════════════════════════════════════════════════════

The alerts page is the supervisor's primary workspace. It presents a
filterable list of alerts with expandable detail for each entry.

\subsection{Filtering and List View}

The supervisor can filter alerts by tier (BLOCK, ALERT, REVIEW), operation
type (retrait, versement, virement, cheque), and review status (pending,
confirmed, rejected). Each alert row displays a color-coded tier badge,
operation type, transaction amount, timestamp, fused anomaly score, and
current review status.

\subsection{Alert Detail}

Expanding an alert reveals the full context needed for a supervisory
decision: client, employee, and branch identifiers; the triggered rules
from the Decision Service with their descriptions; any regulatory flags
with warning icons; and the SHAP explanation.

The SHAP explanation is presented in two layers. The first layer displays
the human-readable narrative generated by the deterministic NLG templates
described in Chapter~\ref{ch:training} --- sentences such as ``This
versement is flagged because the amount is 4.2 times the client's 30-day
average and it is the third deposit within 48 hours.'' The second layer,
hidden behind an expandable section, shows the raw feature-level
contributions: per-feature SHAP values for the Autoencoder and
expected-versus-actual deviations for the LSTM. This two-layer design
serves both audiences: supervisors who need a quick justification and
analysts who want to understand the model's internal reasoning.

\subsection{Supervisor Feedback}

Each pending alert presents two buttons: \textbf{Confirm} and
\textbf{Reject}, with an optional comment field. Clicking either button
sends a \texttt{POST} request to the \texttt{/api/v1/feedback} endpoint
on the FastAPI simulation API, which updates the \texttt{supervisor\_decision}
and \texttt{decision\_timestamp} columns in the \texttt{alerts} table and
upserts a corresponding record into the \texttt{evaluation\_labels} table.
The page then clears its cache and re-renders immediately, reflecting the
updated status without waiting for the next autorefresh cycle.


% ═══════════════════════════════════════════════════════════
\section{KPIs Page}
\label{sec:dashboard-kpis}
% ═══════════════════════════════════════════════════════════

The KPIs page provides an aggregated view of system activity for
operational oversight. Four metric cards at the top display the total
alert count and the count for each of the three actionable tiers (BLOCK,
ALERT, REVIEW). Below the cards, four Plotly charts visualize different
dimensions of the alert data:

\begin{itemize}
    \item A \textbf{pie chart} showing the distribution of alerts across
          tiers, giving an immediate sense of how frequently the system
          escalates to high-severity classifications.
    \item A \textbf{bar chart} grouped by operation type, revealing which
          of the four caisse operations generates the most alerts.
    \item A \textbf{line chart} of alerts per day over the simulation
          period, useful for identifying temporal patterns or spikes.
    \item A \textbf{bar chart} of the top 20 branches by alert count,
          highlighting branches that may warrant closer supervisory
          attention.
\end{itemize}

All aggregation queries use a 30-second cache TTL, trading a small
amount of staleness for reduced database load when multiple supervisors
access the page concurrently.


% ═══════════════════════════════════════════════════════════
\section{Simulation Page}
\label{sec:dashboard-simulation}
% ═══════════════════════════════════════════════════════════

The simulation page allows a supervisor to manually construct a
hypothetical transaction and submit it for scoring without injecting it
into the live pipeline. This serves two purposes: testing whether a
suspicious transaction would trigger an alert before it is processed, and
training supervisors to understand how the system evaluates different
transaction profiles.

The input form collects client ID, employee ID, transaction amount,
operation type, and conditional payload fields (beneficiary ID for
virements, emitter ID and cheque number for cheques). Client-side
validation checks for required fields, positive amounts, and payload
consistency before the form is submitted. The form sends a \texttt{POST}
request to the \texttt{/api/v1/score} endpoint, which runs the full
scoring pipeline synchronously --- enrichment, AE and LSTM inference,
score fusion, SHAP explanation, and decision rules --- and returns the
result.

The response is displayed as a color-coded tier badge, the fused anomaly
score, any triggered rules and regulatory flags, and the full SHAP
explanation in the same two-layer format used on the alerts page.


% ═══════════════════════════════════════════════════════════
\section{Health Page}
\label{sec:dashboard-health}
% ═══════════════════════════════════════════════════════════

The health page queries the Prometheus HTTP API to display the operational
status of the streaming pipeline. Six service status indicators show
whether each service is reachable: the three core streaming services
(enrichment, scoring, decision), the two sink consumers (raw-events sink,
decisions sink), and the simulation API. Each indicator queries the
Prometheus \texttt{up} metric for the corresponding scrape target and
displays a green or red badge.

Below the status indicators, the page displays three operational metrics
derived from the Prometheus instruments described in
Section~\ref{sec:observability}: throughput (events processed per second,
computed from the Counter's rate), p95 processing latency (from the
Histogram), and consumer lag per topic partition (from the Kafka exporter).
If Prometheus is unreachable, the page degrades gracefully to a warning
message rather than crashing, since a monitoring outage should not prevent
the supervisor from accessing other dashboard pages.


% ═══════════════════════════════════════════════════════════
\section{Pipeline Validation and Results}
\label{sec:results}
% ═══════════════════════════════════════════════════════════

With all components integrated --- synthetic data generator, streaming
pipeline, batch profiles, ML models, decision rules, sink consumers, and
dashboard --- the system was validated end-to-end by replaying the full
243,000-event synthetic dataset through the streaming pipeline and
observing the resulting alert distribution.

\subsection{Score Distribution}

Table~\ref{tab:tier-distribution} shows the tier distribution produced by
the validated pipeline after resolving the cold-start and buffer
duplication issues described in Chapter~\ref{ch:streaming}.

\begin{table}[H]
\centering
\caption{Alert tier distribution on the synthetic dataset}
\label{tab:tier-distribution}
\begin{tabular}{@{}lrr@{}}
\toprule
\textbf{Tier} & \textbf{Proportion} & \textbf{Interpretation} \\
\midrule
INFO   & 89\%   & Normal operations, no action required \\
REVIEW & 11\%   & Mild deviations, logged for audit \\
ALERT  & 0.35\% & Significant anomalies, supervisor review \\
BLOCK  & $<$0.1\% & High-confidence anomalies, immediate action \\
\bottomrule
\end{tabular}
\end{table}

The score range spans from 0.015 (routine transactions by established
clients) to 0.999 (high-confidence anomalies matching multiple detection
mechanisms). The distribution is heavily right-skewed, with the vast
majority of events scoring well below the T1 threshold of 0.95. This
shape is expected: in a functioning anomaly detection system, most
transactions should be normal, and the alert tiers should form a steep
pyramid with very few events at the top.

\subsection{Model Performance}

Table~\ref{tab:model-auc} reports the AUC-ROC scores for both models
after retraining on the 30-feature schema.

\begin{table}[H]
\centering
\caption{Model AUC-ROC on the 30-feature schema}
\label{tab:model-auc}
\begin{tabular}{@{}lrr@{}}
\toprule
\textbf{Model} & \textbf{AUC (74 features)} & \textbf{AUC (30 features)} \\
\midrule
Autoencoder & 0.9853 & 0.9714 \\
LSTM        & 0.9866 & 0.9355 \\
\bottomrule
\end{tabular}
\end{table}

The AUC decrease after the feature reduction is expected. The 74-feature
schema included raw profile statistics that made anomalies trivially
separable --- a client whose \texttt{tx\_count\_30d} is zero but who
suddenly transacts is obvious to any model. The 30-feature contrast schema
forces the models to detect anomalies from relative deviations rather than
absolute values, which is harder but more generalizable: the same contrast
features work across archetypes without archetype-specific thresholds.

The LSTM loses more than the Autoencoder because it operates on sequences
of 30-dimensional feature vectors rather than individual events. Each step
in the sequence carries less information per dimension, reducing the LSTM's
ability to model fine-grained temporal patterns. Despite this, both models
remain well above the 0.90 AUC threshold that would indicate reliable
discrimination.

\subsection{Synthetic Data Caveat}

These results validate the system's ability to learn patterns, serve
predictions at streaming latency, and produce a realistic alert
distribution --- but they are not a production performance benchmark. The
synthetic data generator produces anomalies according to known scenarios
with known parameters. Real-world anomalies are messier, more diverse,
and not drawn from the same distribution.

The engineering argument is different from the statistical one: the system
is designed so that plugging in real transaction data requires no code
changes. The generator is replaced by real event sources; the batch
pipeline rebuilds profiles from real history; the models are retrained on
real features computed by the same \texttt{\_compute\_features()} function;
and the retraining triggers (Chapter~\ref{ch:mlops}) ensure that the
models adapt as the data distribution evolves. The AUCs reported here
confirm that this pipeline works end-to-end, not that it will achieve the
same numbers on production data.
